\scanpage{41}
\mspage{39}

que le foncteur $X\unsure{\mapstochar\relbar|\longmapsto}\mathrm{Hom}(X,K)$ transforme le morphisme d'adjonction particulier \unsure{$K$}

\[
 K\xrightarrow{i}\psi\varphi K,
\]

en

\[
 \mathrm{Hom}\bigl(\psi\varphi(K),K\bigr)
 \xrightarrow{u\unsure{\relbar|\mapsto}u\circ i}\mathrm{Hom}(K,K)
\]

qui est une bijection. \deleted{\unsure{Donc}} $\mathrm{id}_K$ est dans l'image, soit $p$ un élément du premier ens qui lui donne naissance

\diagram{\begin{tikzcd}[ampersand replacement=\&]
 K\arrow[r,hook,"i"] \& \psi\underbrace{\varphi(K)}_{\Vert\,K'}\arrow[l,bend left=22,"p"] \&[-1em] \text{donc }p\circ i=\mathrm{id}_K
\end{tikzcd}}

\unsure{Considérons|considérons}

\[
 ip=q:\psi(K')\longrightarrow\psi(K')\qquad(\mathrm{NB}\ q^2=q)
\]

comme $\psi$ est pl. fid. il existe un unique

\[
 q_0:K'\longrightarrow K'\qquad(K'\overset{\mathrm{def}}{=}\varphi(K))
\]

tel que l'on ait

\[
 q=\psi(q_0).
\]

La relation $pi=\mathrm{id}_K$ donne, en appliquant $\varphi$

\[
 \varphi(p)\varphi(i)\inserted{\unsure{\mathrm{fid.}}}=\mathrm{id},\quad\unsure{\text{or}}
\]

donc comme $\varphi(i)$ est iso, (c'est \unsure{évident comm. du|en faisant comm. le} triangle ci-contre),

\diagram{\begin{tikzcd}[ampersand replacement=\&]
 \varphi(K)\arrow[r,"\sim"]\arrow[dr,dash,swap,"{\sim\,\mathrm{id}}"] \& \varphi(\psi\varphi(K))=(\varphi\psi)(\varphi K)\arrow[d,"{\wr\ \mathrm{adj.}}"]\\
 \&\varphi K
\end{tikzcd}}

$\varphi(p)$ l'est. Donc $\varphi(q)=\varphi(i)\varphi(p)$ l'est, puisque $\varphi(i)$ et $\varphi(i)$ le sont. Mais

\[
 \varphi(q)=\varphi\psi(q_0)\xrightarrow[\ ]{\substack{\mathrm{adj}\ \varphi\psi\simeq\mathrm{id}\\ \sim}}q_0,
\]

donc $q_0$ est un iso, donc aussi $q=\psi q_0$. Comme $q^2=q$, cela implique $q=\mathrm{id}_{\unsure{K|{?}}}$. Donc $i$ et $p$ sont des iso. inv. l'un de l'autre, ce qui établit \unsure{(ii)$''$|(ii$'$)} et achève la dém.

\uline{Corollaire} : Considérons deux foncteurs $C^{\unsure{o|0}}\xrightarrow{H}\mathrm{Ens}$, et $C'{}^o\xrightarrow{H'}\mathrm{Ens}$, tels que $H^{\deleted{\unsure{\prime}}}\simeq H'\circ\varphi$, \deleted{\unsure{alors}}. Pour que $H$ soit représentable, il faut et \unsure{s} que $H'$ le soit, et dans ce cas, l'objet $K$ de $C$

% archive sheet number: 40
